% ============ 02. Contexto y guía de lectura ============
\section{Contexto y guía de lectura}

\subsection{Qué es este documento}

Este documento desarrolla el \textbf{criterio 3 de la rúbrica}:
\emph{Documentación y crítica por archivo de código}. El repositorio auditado es el
Observatorio de investigadores reconocidos del Ministerio de Ciencia,
Tecnología e Innovación de Colombia (MinCiencias), que consolida y analiza
\textbf{6 convocatorias (2013--2021)} del padrón oficial de investigadores
y su producción científica. Esta serie cubre los \textbf{6 criterios} de la
rúbrica en documentos separados, más el documento maestro que los integra.

\textbf{En resumen:} Este criterio evalúa la documentación de los archivos de código: qué hace cada uno, cómo está construido, si está bien hecho y si al ejecutarse produce lo que se buscaba. Conclusión: el código está bien construido en general (modular y documentado); 10 de los 16 scripts principales se ejecutan correctamente en un clon del repositorio y los 6 restantes fallan por causas identificadas (datos no versionados y una dependencia no declarada).




\begin{tcolorbox}[nota]
\textbf{Cómo leer este documento:} cada PDF de esta serie desarrolla \emph{un}
criterio de la rúbrica. Los demás criterios se tratan a fondo en sus propios
documentos y en el \textbf{documento maestro} (\texttt{maestro\_auditoria.pdf}).
Si este es tu primer acercamiento, lee primero la portada y este contexto; al
final encontrarás las conclusiones del criterio.
\end{tcolorbox}

\subsection{Glosario (solo los términos que usa este documento)}

Esta tabla incluye únicamente los términos técnicos que aparecen en este
documento, explicados en lenguaje sencillo:

\begingroup
\small
\setlength{\tabcolsep}{3pt}
\begin{longtable}{>{\raggedright\arraybackslash}p{4.3cm}>{\raggedright\arraybackslash}p{9.8cm}}
\caption{Glosario de términos presentes en este documento}
\label{tab:glosario}\\
\toprule
\textbf{Término} & \textbf{Qué significa} \\
\midrule
\endfirsthead
\toprule
\textbf{Término} & \textbf{Qué significa} \\
\midrule
\endhead
    HHI (Herfindahl-Hirschman) & Índice que mide si algo está concentrado o repartido. Va de 0 (todo repartido) a 1 (todo en un solo lugar). En el informe se usa para ver si los investigadores se concentran en pocos departamentos.
 \\
    Matriz de transición & Tabla que muestra con qué probabilidad un investigador pasa de una categoría (p. ej. Junior) a otra (p. ej. Asociado) entre dos convocatorias. Permite ver si el sistema 'sube', 'mantiene' o 'pierde' investigadores.
 \\
    Tasa de retención & Porcentaje de investigadores que siguen presentes en la siguiente convocatoria. Una retención baja puede indicar abandono o cambio de reglas.
 \\
    Representación relativa & Cuántas veces aparece un grupo (p. ej. afrocolombianos) en el padrón frente a lo que correspondería según su peso en la población colombiana (DANE). Un valor de 0.3 indica que aparece 3 veces menos de lo esperado.
 \\
    k-prototypes & Algoritmo de agrupamiento (clustering) que agrupa personas con perfiles similares cuando los datos mezclan números (edad) y categorías (área, género).
 \\
    MCA / FAMD & Técnicas para resumir tablas con muchas categorías en pocos ejes visualizables, de modo que se puedan ver patrones (p. ej. qué áreas se asocian con qué géneros).
 \\
    Modelo dimensional (esquema estrella) & Forma de organizar los datos en tablas de 'hechos' (métricas: quién, cuándo, qué categoría) y 'dimensiones' (descripciones: área, región, institución), para hacer consultas rápidas y reportes.
 \\
    DuckDB & Base de datos analítica ligera que se ejecuta dentro del propio proyecto (sin instalar un servidor). Se usa para almacenar el modelo dimensional.
 \\
    Mojibake & Texto ilegible por un problema de codificación (p. ej. 'Física' en lugar de 'Física'). Detectado en el CSV del repositorio.
 \\
    Contrato de esquema & Validación automática que garantiza que un archivo de datos tenga las columnas, tipos y valores esperados antes de analizarlo.
 \\
    Pipeline & Cadena de procesos que lleva los datos desde la fuente hasta el análisis (ingesta, limpieza, transformación, modelado). También llamada 'cadena de procesamiento' o 'flujo de datos'.
 \\
    Markovianidad & Propiedad de un proceso en el que la probabilidad de pasar al siguiente estado solo depende del estado actual (y no de toda la historia). Se usa para validar si las matrices de transición de categorías son adecuadas.
 \\
    Modelos de supervivencia (Kaplan-Meier / Cox) & Técnicas estadísticas para analizar 'tiempo hasta que ocurre un evento' (p. ej. cuánto tarda un investigador en salir del sistema), incluyendo datos incompletos.
 \\
    Intervalo de confianza (IC) y bootstrap & El intervalo de confianza indica el rango plausible de una estimación (p. ej. 0.15-0.19). El bootstrap es un método que lo calcula remuestreando los datos muchas veces, sin fórmulas complejas.
 \\
    CI/CD y Docker & Automatización de pruebas y despliegue (CI/CD) y empaquetado de la aplicación en un contenedor (Docker) para que funcione igual en cualquier máquina.
 \\
    Sprint & Ciclo corto de trabajo (en este proyecto, de 2 semanas) con objetivos concretos.
 \\
\bottomrule
\end{longtable}
\endgroup
